\section{Diagnosing the sharp remainder without fitting its main term}
\label{sec:sharpdiagnostics}

The alternatives for the sharp primitive remainder are not exhausted by
$L$ and $L\log L$. Strict convexity already gives the qualitative
improvement $o(L\log L)$, proved below, but supplies no effective
intermediate rate. A normalized M\"obius formula identifies a sufficient
criterion for $O(L)$. An exact dyadic inversion removes the leading area
constant from the numerical test, and an inverse-flatness formula measures
the contribution of a particular supporting row. Finally we describe the
dyadic diagnostic using hyperbolic trace recursion on two tori, distinguish
the supplied draft's unreplayed numerical report from a new bounded
replacement calculation, and provide the latter's computational recipe.
The proofs and the exploratory numerical observations are stated separately.

\subsection{The normalized all-point discrepancy and the logarithmic loss}
\label{subsec:mobiusmodulus}

Rescale the norm so that its shortest nonzero integer vector has length one.
Write $a=\operatorname{Area}(B)$,
\[
 A_0(u)=\#(uB\cap(\mathbb Z^2\setminus\{0\})),\qquad
 D_B(u)=A_0(u)-a u^2,
\]
and let $P_B(L)$ count primitive vectors. For $M=\lfloor L\rfloor$,
M\"obius inversion gives the exact identity
\begin{equation}
 \frac{P_B(L)-aL^2/\zeta(2)}{L}
 =\sum_{d\le M}\frac{\mu(d)}d\,
       \frac{D_B(L/d)}{L/d}
   -aL\sum_{d>M}\frac{\mu(d)}{d^2}.
 \label{eq:primitiveexactnormalizedmobius}
\end{equation}
The final term is bounded uniformly for $L\ge2$. Thus the linear
remainder problem is precisely a boundedness problem for the displayed
M\"obius-weighted harmonic sum. The prime number theorem alone does not
bound it: the geometric weight $D_B(L/d)/(L/d)$ varies with both $L$ and $d$.

For a convenient sufficient condition, put
\[
 \omega_B(u)=\sup_{v\ge u}\frac{|D_B(v)|}{v},\qquad u\ge1.
\]
The elementary boundary estimate makes $\omega_B$ bounded, and it is
nonincreasing. Absolute values in~\eqref{eq:primitiveexactnormalizedmobius}
give
\begin{equation}
 |P_B(L)-aL^2/\zeta(2)|
 \le C_B L\left(1+\int_1^L\omega_B(u)\,\frac{du}{u}\right).
 \label{eq:primitiveDini}
\end{equation}
Indeed, for $d\le M$,
\[
 \int_{L/(d+1)}^{L/d}\omega_B(u)\frac{du}{u}
 \ge\omega_B(L/d)\log(1+1/d)
 \ge\frac{\omega_B(L/d)}{2d},
\]
with the harmless interval below $1$ treated separately.
Consequently the Dini condition
\[
 \int_1^\infty\omega_B(u)\,\frac{du}{u}<\infty
\]
implies an $O_B(L)$ primitive remainder without any M\"obius cancellation.
With the prime number theorem, this improves to $o_B(L)$; see
\eqref{eq:primitiveDiniPNT} below.
In particular any all-point power saving $D_B(u)=O_B(u^{1-\eta})$,
$\eta>0$, suffices. Infinite graph-flatness does not logically exclude
this route: for example a bound $u\exp(-\sqrt{\log u})$ has no fixed
power saving but has an integrable normalized modulus.

\begin{proposition}[A qualitative strict-convexity improvement]
\label{prop:strictconvexlittleo}
For every fixed strictly convex compact body $B\subset\mathbb R^2$
containing the origin in its interior,
\[
 D_B(L)=o_B(L),\qquad
 P_B(L)-\frac{a}{\zeta(2)}L^2=o_B(L\log L).
\]
These statements supply no quantitative rate of little-oh convergence.
\end{proposition}

\begin{proof}
Fix a unit vector $\nu$ and let $w_\nu(s)$ be the length of the
chord $B\cap\{x:x\cdot\nu=s\}$. Its support is an interval
$[a_\nu,b_\nu]$. Convexity makes $w_\nu$ concave, and strict
convexity makes both extreme supporting chords singletons, so
$w_\nu(a_\nu)=w_\nu(b_\nu)=0$. A finite continuous concave
function on a compact interval is absolutely continuous: its
monotone one-sided derivative is integrable, and the fundamental
theorem on interior subintervals extends to the endpoints.
The zero extension therefore belongs to $W^{1,1}(\mathbb R)$.
Moreover
\[
 \|w_\nu'\|_1=2\max_s w_\nu(s)\le2\operatorname{diam}(B).
\]
Fubini and integration by parts give
\[
 \widehat{\mathbf1_B}(r\nu)=\widehat w_\nu(r),\qquad
 r\widehat{\mathbf1_B}(r\nu)
   =\frac1{2\pi i}\widehat{w_\nu'}(r)\longrightarrow0
\]
by the Riemann--Lebesgue lemma. The same formula gives the uniform
bound $|\widehat{\mathbf1_B}(\xi)|\le C_B/|\xi|$.

Fix a nonnegative smooth function $\rho_\epsilon$ of integral one,
supported in the Euclidean disk of radius $\epsilon$. Choose $r_0>0$
with $r_0\mathbb D\subset B$, and put $\delta=\epsilon/r_0$.
Convexity gives the pointwise sandwich
\[
 \mathbf1_{(L-\delta)B}*\rho_\epsilon
 \le\mathbf1_{LB}
 \le\mathbf1_{(L+\delta)B}*\rho_\epsilon.
\]
For this fixed $\epsilon$, Poisson summation shows that either
smoothed lattice count equals its area plus $o_{B,\epsilon}(L)$: after division
by $L$, each nonzero-frequency term tends to zero, and the summable
majorant is a constant times
$|m|^{-1}|\widehat\rho_\epsilon(m)|$.
The two areas differ from $aL^2$ by $O_B(\epsilon L)+O_B(\epsilon^2)$.
Thus
\[
 \limsup_{L\to\infty}|D_B(L)|/L\le C_B\epsilon.
\]
Letting $\epsilon\downarrow0$ proves the first claim. It implies
$\omega_B(u)\to0$, and logarithmic Ces\`aro averaging in
\eqref{eq:primitiveDini} proves the second.
\end{proof}

The qualitative improvement excludes a fixed positive multiple of
$L\log L$ along a subsequence, but still allows $L\log\log L$,
$L(\log L)^{0.9}$, and many other intermediate scales. Quantitative
control of $\omega_B$, or signed cancellation between its M\"obius
blocks, is needed to reach $O_B(L)$.

\subsection{A rational-polygon comparison}
\label{subsec:rationalpolygonbenchmark}

The intermediate scales cannot be dismissed merely because the
$O(L\log L)$ proof takes absolute values. For a rational polygon
$\mathcal A$ that is star-shaped about the origin, B\'ar\'any,
Martin, Naslund, and Robins prove
\[
 E_{\mathcal A}(L)
 =\Omega_\pm\bigl(L\sqrt{\log\log L}\bigr),
\]
while their upper bound for every rational polygon is
\[
 E_{\mathcal A}(L)
 =O_{\mathcal A}\bigl(L(\log L)^{2/3}
                         (\log\log L)^{4/3}\bigr).
\]
These are Theorems~1.4 and~1.2, respectively, in
\emph{Primitive points in rational polygons}~\cite{BaranyPrimitivePolygons}.
The polygons have genuine straight sides. These results do not
transfer to a strictly convex McShane--Rivin norm ball; they show
why an unspecified arithmetic cancellation is not a proof of a
linear error. Even the elementary square illustrates the relevant
arithmetic: at integer $N\ge1$,
\[
 \#([-N,N]^2\cap\mathbb Z^2_{\rm prim})
 =8\sum_{n=1}^N\varphi(n).
\]
There is also a numerical warning. Freezing a rational polygonal
approximation to a strictly convex ball can introduce the polygon's own
totient fluctuations. A fixed area error $\delta a$ contributes
$\delta a L^2$; resolving a linear remainder requires area accuracy
$o(L^{-1})$, together with control of points near the two boundaries.
The trace enumeration below avoids replacing the norm ball by a polygon.

\subsection{An exact test independent of the leading coefficient}
\label{subsec:dyadicdiagnostic}

Suppose $N(L)/L^2\to c$, put $E(L)=N(L)-cL^2$, and define
\[
 \Delta_2N(L)=N(2L)-4N(L)=E(2L)-4E(L).
\]
Iterating the latter equality and using $E(L)=o(L^2)$ gives
\begin{equation}
 E(L)=-\sum_{j=0}^\infty4^{-j-1}\Delta_2N(2^jL).
 \label{eq:dyadicerrorinverse}
\end{equation}
The series is understood as the limit of its telescoping partial
sums; it is absolutely convergent under any of the bounds below.
Consequently
\[
 E(L)=O(L)\quad\Longleftrightarrow\quad
 \Delta_2N(L)=O(L),
\]
and likewise
\[
 E(L)=O(L\log L)\quad\Longleftrightarrow\quad
 \Delta_2N(L)=O(L\log L).
\]
For example, if $|\Delta_2N(u)|\le Cu$ for every $u\ge L$, then
\[
 |E(L)|\le\frac{CL}{4}\sum_{j\ge0}2^{-j}=\frac{CL}{2}.
\]
These are asymptotic equivalences, not finite-computation certificates.
They eliminate contamination by an inaccurately estimated value of $c$.

If an estimated coefficient $\widetilde c=c+\delta c$ is subtracted
instead, then
\[
 \frac{N(L)-\widetilde cL^2}{L}
 =\frac{E(L)}L-\delta cL.
\]
Thus a small coefficient error creates a growing normalized drift. The
dyadic statistic removes this source of error exactly. It does not
remove the possibility of rare large excursions beyond the computed range.
For measurements over a block define
\begin{align}
 M(T)&=\sup_{T\le L\le2T}\frac{|\Delta_2N(L)|}{L},\notag\\
 Q(T)&=\left(\frac1T\int_T^{2T}
                    \left|\frac{\Delta_2N(L)}L\right|^2\dd L\right)^{1/2}.
 \label{eq:dyadicblockmetrics}
\end{align}
The first detects the largest computed excursions, while the second
measures their contribution over the whole block. Between events at
$L=\ell_\gamma$ and $L=\ell_\gamma/2$, $\Delta_2N(L)$ is constant.
Hence its normalized absolute value decreases between events, and
the block supremum is determined by both one-sided event limits and the
block endpoints. All events at the same location must be applied together;
an arbitrary order of simultaneous positive and negative jumps can create
spurious intermediate values. The left limit at the initial endpoint
$T$ lies outside the block and is excluded.
Its squared integral over such an interval $[a,b]$ is exactly
$D^2(1/a-1/b)$. A coarse integer mesh can miss narrow excursions.


\subsection{A quantitative row-crossing formula}
\label{subsec:inverseflatness}

The exponent-one obstruction does not specify the slowly varying factor
in the discrepancy. The local graph itself gives a more precise statement.
Let $B\subset\RR^2$ be strictly convex, compact, and contain the origin in
its interior. After an integral unimodular change of coordinates, suppose
its unique top point is $p=(\alpha,h)$, with $h>0$ and
$\alpha/h\notin\mathbb Q$. Write the upper boundary near $p$ as
\[
 y=h-f(x-\alpha),\qquad f(0)=0.
\]
For sufficiently small $u\geq0$, define the one-sided inverse widths
\[
\begin{gathered}
 \rho_+(u)=\sup\{s\geq0:f(s)\leq u\},\qquad
 \rho_-(u)=\sup\{s\geq0:f(-s)\leq u\},\\
 W(u)=\rho_+(u)+\rho_-(u),
\end{gathered}
\]
where the suprema are restricted to the local graph interval. Strict
convexity implies $W(u)\to0$ as $u\downarrow0$. The functions
$\rho_\pm$, and hence $W$, are concave on a sufficiently small interval:
they are inverses of increasing convex functions vanishing at zero.
No symmetry between the two sides of the graph is assumed.

Let
\[
 P_B(L)=\#(LB\cap\ZZ^2_{\rm prim}),\qquad
 c_B=\frac{\operatorname{Area}(B)}{\zeta(2)},\qquad
 E_B(L)=P_B(L)-c_BL^2.
\]

\begin{proposition}[Inverse flatness controls a primitive-row burst]
\label{prop:inverseflatrow}
Let $n$ tend to infinity through primes, put $L=n/h$, and let
$\delta=\delta(L)>0$ satisfy $\delta=o(L)$. Set $T=L+\delta$.
For all sufficiently large $n$, the intersection of $TB$ with the row
$y=n$ is the interval with endpoints
\[
 a_T=T\alpha-T\rho_-\!\left(\frac{h\delta}{T}\right),\qquad
 b_T=T\alpha+T\rho_+\!\left(\frac{h\delta}{T}\right).
\]
Its width is
\begin{equation}
 w_n(\delta)=T W\!\left(\frac{h\delta}{T}\right).
 \label{eq:inverseflatwidth}
\end{equation}
The number of primitive lattice points on this row is $w_n(\delta)+O(1)$,
with an absolute rounding constant, and
\begin{equation}
 \max\{|E_B(L)|,|E_B(L+\delta)|\}
 \geq\frac12\left[
 w_n(\delta)-c_B(2L\delta+\delta^2)-2\right].
 \label{eq:inverseflatdiscrepancy}
\end{equation}
In particular, with $\delta=L^{-1}$,
\begin{equation}
 \max\{|E_B(L)|,|E_B(L+L^{-1})|\}
 \geq\frac L2W(h/L^2)-O_B(1).
 \label{eq:inverseflatmicroscopic}
\end{equation}
\end{proposition}

\begin{proof}
At scale $L=n/h$, the row touches $LB$ only at $(L\alpha,n)$.
The first coordinate is irrational, so there is no lattice point there.
At scale $T$, the normalized row has height $n/T$ and thus depth
\[
 h-\frac nT=\frac{h\delta}{T}
\]
below the supporting line. Inverting the two sides of the graph gives
the displayed endpoints. Since $\delta/T\to0$ and $W(u)\to0$,
\[
 \frac{w_n(\delta)}n=\frac{T}{hL}
 W\!\left(\frac{h\delta}{T}\right)\longrightarrow0.
\]
Eventually $w_n(\delta)<n$. The number of integers in a closed interval
of width $w$ differs from $w$ by at most one. Such an interval contains
at most one multiple of the prime $n$. Every other integer $k$ gives a
primitive pair $(k,n)$. Therefore the row contains at least $w-2$ and
at most $w+1$ primitive lattice points.

Since $LB\subset TB$, all these points are new, while points elsewhere
cannot disappear. Thus
\[
 E_B(T)-E_B(L)
 \geq w_n(\delta)-2-c_B(T^2-L^2),
\]
which implies~\eqref{eq:inverseflatdiscrepancy}.

For $\delta=L^{-1}$, set $u_0=h/L^2$. Concavity and $W(0)=0$ give
\[
 W\!\left(\frac h{L^2+1}\right)
 \geq\frac{L^2}{L^2+1}W(u_0).
\]
Multiplying by $T=(L^2+1)/L$ yields
$w_n(L^{-1})\geq LW(h/L^2)$. The change in the quadratic main term
is $c_B(2+L^{-2})$, proving~\eqref{eq:inverseflatmicroscopic}.
\end{proof}

The theorem remains useful without a model formula for $f$. It measures
the width of the graph at depth $L^{-2}$ and converts that width into a
lower bound for the discrepancy. Infinite flatness implies, for every
$N$, a lower bound $W(u)\gg_N u^{1/N}$ and recovers
Proposition~\ref{prop:normballflatrow}. A specified modulus yields a
specified slowly varying factor.


% BEGIN INLINE FIGURE: 07_flat_row.tex
\begin{figure}[htbp]
\centering
\begin{tikzpicture}[font=\small,>=Latex,line cap=round]
  \begin{scope}[shift={(-3.7,.3)}]
    \node[font=\bfseries\small] at (0,2.25) {Local inverse width};
    \draw[mpgray!60,dashed] (-2.05,1.5)--(2.05,1.5);
    \draw[mpblue,line width=1.1pt,domain=-1.85:1.85,samples=100]
      plot (\x,{1.5-.18*\x*\x*\x*\x});
    \draw[mpteal,line width=.8pt] (-2.05,.8)--(2.05,.8);
    \fill[mpblue] (0,1.5) circle (2pt) node[above=3pt] {$p=(\alpha,h)$};
    \fill[mpteal] (-1.404,.8) circle (1.8pt);
    \fill[mpteal] (1.404,.8) circle (1.8pt);
    \draw[mpgray!55,dashed] (-1.404,.8)--(-1.404,.35) (1.404,.8)--(1.404,.35);
    \draw[<->,mpteal] (-1.404,.35)--node[below=2pt] {$W(u)$}(1.404,.35);
    \draw[<->,mpgray] (1.95,.8)--node[right] {$u$}(1.95,1.5);
    \node[left=2pt,text=mpteal] at (-2.05,.8) {$h-u$};
    \node[text=mpblue,align=center,font=\scriptsize] at (0,-.65)
      {$y=h-f(x-\alpha)$\\profile shown schematically};
  \end{scope}
  \begin{scope}[shift={(3.2,1.15)}]
    \node[font=\bfseries\small] at (0,1.4) {A prime lattice row};
    \fill[mpteal!9] (-1.65,-.13) rectangle (1.85,.13);
    \draw[->,mpgray] (-2.65,0)--(2.65,0) node[right] {$y=n$};
    \foreach \x in {-2.5,-2,-1.5,-1,-.5,.5,1,1.5,2,2.5}
      \fill[mpgray!60] (\x,0) circle (1.8pt);
    \foreach \x in {-1.5,-1,-.5,.5,1,1.5}
      \fill[mpteal] (\x,0) circle (2.3pt);
    \filldraw[fill=white,draw=mpred,line width=1pt] (0,0) circle (2.8pt);
    \draw[mpgray] (-1.65,-.2)--(-1.65,.2) (1.85,-.2)--(1.85,.2);
    \node[above=4pt] at (-1.65,0) {$a_T$};\node[above=4pt] at (1.85,0) {$b_T$};
    \draw[<->,mpteal] (-1.65,-.43)--node[below=2pt] {$w_n(\delta)$}(1.85,-.43);
    \node[align=center,text=mpred,font=\scriptsize] at (0,-1.25)
      {At most one excluded point\\when $w_n(\delta)<n$: $k\equiv0\pmod n$};
    \node[font=\scriptsize,text=mpgray] at (0,.8) {Teal points are primitive; $n$ is prime};
  \end{scope}
  \node[align=center] at (0,-1.75)
    {$L=n/h,\qquad T=L+\delta,\qquad
      u=\displaystyle\frac{h\delta}{T},\qquad
      w_n(\delta)=T\,W(u)$.};
\end{tikzpicture}
\caption[A flat supporting cap and a primitive-row burst]{How a supporting cap becomes a primitive-row contribution.
At the contact scale $L=n/h$, the hypothesis $\alpha/h\notin\mathbb Q$
leaves the row empty. Increasing the scale to $T=L+\delta$ cuts an interval
of width $TW(h\delta/T)$ from that row. When $n$ is prime and the width
is smaller than $n$, all but at most one of its integer points are
primitive. The local profile and row spacing are schematic; the
displayed width relation is exact.}
\label{fig:mono-flat-row}
\end{figure}
% END INLINE FIGURE

\subsection{Model flatness profiles and what they do not determine}
\label{subsec:flatnessmodels}

For illustration, suppose the two-sided local graph satisfies
\[
 f(s)\leq C\exp\!\left(-\frac a{|s|^\beta}\right),
 \qquad a,\beta>0,
\]
for sufficiently small nonzero $s$. Then
\[
 W(u)\geq
 2\left(\frac a{\log(C/u)}\right)^{1/\beta},
\]
and~\eqref{eq:inverseflatmicroscopic} gives, along arbitrarily large
prime-row contact scales,
\begin{equation}
 \max\{|E_B(L)|,|E_B(L+L^{-1})|\}
 \gg_{B,a,\beta}\frac L{(\log L)^{1/\beta}}.
 \label{eq:exponentialflatrow}
\end{equation}
If the graph has matching upper and lower estimates of this form, the
row width itself has this order. This does not prove that the full
discrepancy has exactly this order; other rows and other supporting
directions can contribute.

Other infinitely flat local profiles give different slowly varying
factors. For example, the symmetric model
$f(s)=\exp[-\exp(|s|^{-\beta})]$ gives a row width of order
$L/(\log\log L)^{1/\beta}$. The model
\[
 f(s)=\exp[-(\log(1/|s|))^\gamma],\qquad\gamma>1,
\]
gives a width of order
\[
 L\exp[-(2\log L)^{1/\gamma}].
\]
These are convex for sufficiently small $|s|$ and illustrate local
possibilities. They are not asserted here to be the precise profiles
of a hyperbolic torus norm ball.

For a strictly convex body, $W(u)\to0$. Consequently the single-row
argument with a microscopic scale increment yields a quantity $o(L)$,
although it can exceed every fixed power $L^{1-\varepsilon}$.
It neither proves a literal $\Omega(L)$ lower bound nor distinguishes
an $O(L)$ global remainder from an $O(L\log L)$ remainder. Its role is
to quantify one geometric obstruction, rather than to identify the
entire sharp remainder.

\subsection{An exact row diagnostic and its precision requirements}
\label{subsec:rowdiagnostic}

For an arbitrary positive integer $n$ and a row interval $[a,b]$, the
primitive count is exactly
\begin{equation}
 \#\{k\in\ZZ\cap[a,b]:\gcd(k,n)=1\}
 =\sum_{d\mid n}\mu(d)
 \left(\left\lfloor\frac bd\right\rfloor
       -\left\lceil\frac ad\right\rceil+1\right).
 \label{eq:primitiveintervalexact}
\end{equation}
In particular, it equals
\[
 (b-a)\frac{\varphi(n)}n+O(2^{\omega(n)}).
\]
The formula includes endpoints exactly. For prime $n$ and $b-a<n$,
the simpler absolute-error comparison with $b-a$ used above applies.

One can therefore record the local statistic
\[
 J_n=P_{B,n}(n/h+h/n),\qquad
 P_{B,n}(T)=\#\{k\in\ZZ:\gcd(k,n)=1,\ (k,n)\in TB\},
\]
at prime $n$. Since $L=n/h$ and $h/n=L^{-1}$, the proposition gives
\[
 J_n=(L+L^{-1})W\!\left(\frac h{L^2+1}\right)+O(1).
\]
This statistic needs no estimate of $\operatorname{Area}(B)$; it tests
the local flatness mechanism directly. It does require a verified
lattice-rational supporting normal at an irrational radial point.
Finding such a direction is a separate geometric step and is not
automatic for an arbitrary fixed surface.

The height $h$ must also be known accurately. An error $\Delta h$ moves
the contact scale $n/h$ by approximately $-L\Delta h/h$. Resolving
the window of width $L^{-1}$ thus requires $|\Delta h|\ll hL^{-2}$.
Using a larger window relaxes this requirement but changes the depth
sampled by~\eqref{eq:inverseflatwidth} and increases the background
main-term increment. The local statistic is consequently complementary
to a global main-term-free discrepancy diagnostic; it does not replace
one.



\subsection{A trace diagnostic: retained report and bounded replacement}
\label{subsec:torustraceexperiment}

The count here is $N_X(L)$ for unoriented essential simple closed
geodesics on a once-punctured torus, with the convention
$N_X(L)=\#\{\gamma:\ell_\gamma\le L\}$. Each primitive pair
$(p,q)\in\mathbb Z^2$ is identified with $(-p,-q)$, so this is half
the signed primitive-vector count for the corresponding norm ball.
A positive trace triple satisfies
\[
 x^2+y^2+z^2=xyz,\qquad x,y,z>2,
\]
and a Farey edge replaces $z$ by $xy-z$. The identification of Farey
regions with simple curves and this edge relation are described in
\cite[Sections~1--2]{BowditchMarkoff}. For the experiment take the reduced
triples
\begin{equation}
 (x,y,z)=(3,3,3),\qquad
 (x,y,z)=\left(\frac52,\frac{10-\sqrt{11}}2,4\right).
 \label{eq:diagnostictracetriples}
\end{equation}
The first is the modular torus; the second is an asymmetric example.
Both satisfy $2<x\le y\le z\le xy/2$.

Assign the root curves the slopes $(1,0)$, $(0,1)$, and $(1,1)$,
respectively. The three open sectors between these roots enumerate
every other unoriented primitive slope once by Farey mediants. In the
negative-slope sector use the signed endpoints $(1,0),(0,-1)$;
the other two sectors have endpoints $(1,0),(1,1)$ and
$(0,1),(1,1)$. Canonical representatives have $q>0$, with $(1,0)$
as the only representative at infinity. Sector endpoints are counted
only as roots. After introducing an outward
trace $c\ge\max(a,b)>2$, the two new traces are $ac-b$ and $bc-a$,
both strictly larger than $c$. Therefore pruning a sector when its new
curve exceeds the cutoff does not omit shorter descendants. Equal traces
are retained with their curve multiplicities; they are not deduplicated.
The hyperbolic length associated with trace $x$ is $2\arcosh(x/2)$.

To avoid overflow the computation uses logarithmic traces. If
$u_a=\log a$, $u_b=\log b$, and $u_c=\log c$, then
\[
 \log(ab-c)=u_a+u_b+
        \log\bigl(1-\exp(u_c-u_a-u_b)\bigr).
\]
On the chosen outward edges the subtracted ratio is at most $1/2$,
so this computation avoids catastrophic cancellation. The supplied draft
describes extended-precision arithmetic for this recursion and double
precision for sorting the resulting lengths. That reported implementation
was not included in the supplied materials. The bounded replacement below
uses binary64 logarithmic recursion and a separate exact-trace oracle.
Neither calculation provides an interval certificate of every event location.

\paragraph{Retained numerical report, awaiting reproduction.}
The supplied draft reports enumeration through hyperbolic length $8192$, obtaining
$18\,191\,508$ curves for the modular torus and $18\,486\,116$ for
the asymmetric torus. Such a cutoff permits evaluating the dyadic block statistics
through $T=2048$, because the largest count used in a block is $N_X(4T)$.
The following selected values are retained verbatim from that report;
all entries are rounded and have not been reproduced for this edition.

\begin{center}
\begin{tabular}{r r r r r}
\toprule
&\multicolumn{2}{c}{Modular torus}&\multicolumn{2}{c}{Asymmetric torus}\\
$T$&$M(T)$&$Q(T)$&$M(T)$&$Q(T)$\\
\midrule
32 & 2.1686 & 0.6015 & 0.8134 & 0.2734 \\
128 & 1.1130 & 0.3501 & 0.5285 & 0.1494 \\
512 & 0.7415 & 0.1697 & 0.3756 & 0.0950 \\
2048 & 0.4342 & 0.0950 & 0.2492 & 0.0612 \\
\bottomrule
\end{tabular}
\end{center}

According to the supplied draft, the reported maxima use all computed
length events, including the half-length events, rather than an integer
mesh. It reports that an independent direct-trace computation at $100$
decimal digits reproduces
all lengths through $64$ within $3.6\times10^{-15}$, and thirty
selected long recursion paths reaching lengths near $8192$ agree within
$7.5\times10^{-13}$. The latter include long repeated-twist paths.
The same report states that every dyadic cutoff used for the integer counts is separated from the
computed spectrum by at least $8.3\times10^{-5}$.

It also describes a separate sensitivity calculation treating nearby events as having an
unknown internal order and allows errors up to $10^{-6}$ in each event
position. Over all blocks the resulting lower and upper bounds for
$M(T)$ differ by less than $7.1\times10^{-7}$. This is a conditional
sensitivity check, not a proof that every computed length has that error
bound. The advertised source, full block table, precision report, and
sensitivity tables were not found in the supplied archive or repositories.
Consequently these counts, table entries, precision comparisons, and
sensitivity bounds remain unreplayed reported results. This edition does
not claim that the advertised diagnostic archive accompanies the book.

The displayed finite-range normalized excursions decrease overall;
they do not establish an asymptotic linear or sublinear bound. The modular
torus was already singled out as a candidate for exceptional behavior in
\cite[Conjecture~4.3]{McShaneRivinTori}. Neither example is asserted to
have a verified lattice-rational supporting normal at an infinitely flat
irrational point. Sparse resonances can occur beyond the tested range.
The stability claims for this reported long run remain awaiting their
computational evidence.

\paragraph{A bounded replacement calculation.}
The maintained recipe \path{experiments/section30/run.py} performs a new
calculation through length $64$ for the same two exact trace triples.
It finds $1104$ modular and $1127$ asymmetric curves. This is a replacement
experiment, not a reconstruction of the missing length-$8192$ run.
Its complete slope list, exact traces, numerical lengths, and block
statistics are recorded in \path{artifacts/section30/result.json}, with
the source identity and precision settings. The new block values are
\begin{center}
\begin{tabular}{r r r r r}
\toprule
&\multicolumn{2}{c}{Modular torus}&\multicolumn{2}{c}{Asymmetric torus}\\
$T$&$M(T)$&$Q(T)$&$M(T)$&$Q(T)$\\
\midrule
4 & 3.3289 & 1.4538 & 1.4983 & 0.7563 \\
8 & 3.2433 & 1.2476 & 1.2352 & 0.4951 \\
16 & 2.0238 & 0.8521 & 1.2246 & 0.3942 \\
\bottomrule
\end{tabular}
\end{center}

The check uses two enumerations. The first traverses the three sectors
with logarithmic traces. The second evaluates exact trace polynomials
for all primitive slopes in the box $|p|\le160$, $0\le q\le160$,
using Euclidean Farey parents, without length-based pruning. This box
suffices for these two fixtures: put $v=\log(\operatorname{tr}/2)$.
The outward subtraction bound gives $v_c\ge v_a+v_b$, and each root
has $v\ge\log(5/4)>1/5$. The corresponding sector weight is
$F=|p|+q$ for $p<0$, $F=p$ for $p\ge q\ge0$, and $F=q$ for
$q\ge p\ge0$. Thus $F\ge\max(|p|,q)$ and
\[
 \ell=2\arcosh(e^v)>2v\ge\frac25F.
\]
Every curve of length at most $64$ therefore lies in the stated box.
There are $31\,224$ primitive unoriented slopes in that box, including
$(1,0)$. Equal traces retain all their distinct slopes.

The oracle uses rational arithmetic in $\mathbb Q(\sqrt{11})$ for
traces and evaluates lengths at $100$ and $140$ decimal digits.
Opposite-signed algebraic coefficients are evaluated through the exact
norm and the conjugate denominator to avoid subtractive cancellation.
Both precisions select the same slopes as the logarithmic traversal;
the length comparisons differ by less than $1.1\times10^{-98}$ between
the two decimal precisions and by less than $1.1\times10^{-13}$ from
the binary64 logarithmic calculation. These are observed residuals,
not rigorous error enclosures.

The displayed metrics use the oracle lengths rounded to binary64.
An event sweep is compared with direct counts on every event interval,
including one-sided limits and simultaneous jumps. Comparing metrics
from the logarithmic and exact-trace length calculations gives differences
below $2.7\times10^{-15}$ for $M$ and $3.2\times10^{-15}$ for $Q$.
The recipe also supplies conservative rational bounds for $M$, conditional
on each stored length having error at most $10^{-10}$; it does not
certify that error hypothesis. These bounds concern the fixed supplied
curve list; application to the complete spectrum also assumes that no
omitted curve can enter a block threshold under the perturbation.
The full derivation and run instructions
are in \path{experiments/section30/README.md}. This bounded check provides
no long-range sensitivity bound, no asymptotic remainder theorem, and
no audit of the remaining proof chain in this section.

\subsection{A recent qualification of irrational flatness}
\label{subsec:recentflatnessqualification}

Doan--Li--Nguyen's recent preprint~\cite{DoanLiNguyenStableNorm},
version~3 dated 9 October 2026, gives a Diophantine refinement of the
irrational flatness statement. Its Theorem~4.7 states that infinite
graph-flatness holds precisely when the exponential rational-approximation
exponent vanishes; superexponentially approximable irrational directions
can have graph-flatness order only one. In slope coordinates the relevant
quantity is $\limsup_j(\log q_{j+1})/q_j$, where $q_j$ are continued-fraction
denominators. The local convexity mechanism is that a nearby rational
corner can remain visible at a much smaller spatial scale than its
exponentially small angle.

We have checked the local Fricke and convexity mechanism relevant to
this distinction, but have not audited the entire preprint. None of its
global counting assertions is used in the proofs above. In particular,
Propositions~\ref{prop:normballflatrow} and~\ref{prop:inverseflatrow}
retain their explicit flat-point hypotheses, and
Proposition~\ref{prop:strictconvexlittleo} uses only strict convexity.
The preprint also continues to list the McShane--Rivin $O_X(L\log L)$
bound as the cumulative counting benchmark; its multiplicity estimate
concerns exact length levels, a different question.


